\begin{table}[H]
\centering
\caption{Nama Daerah Perbandingan Cluster Statis vs IPM Rata-rata (UNDP)} \label{tab:samplestatis_ipm_avg_name}
\begin{adjustbox}{max width=\textwidth}
\begin{tabular}{|p{4cm}|p{10cm}|}
\hline
\textbf{Kategori} & \textbf{Daftar Wilayah} \\
\hline
Kl. 1 (Rendah) \& Sangat Tinggi & Kota Jakarta Pusat \\
\hline
Kl. 1 (Rendah) \& Tinggi & Kota Cilegon \\
\hline
Kl. 1 (Rendah) \& Sedang & - \\
\hline
Kl. 1 (Rendah) \& Rendah & - \\
\hline
Kl. 2 (Sedang) \& Sangat Tinggi & - \\
\hline
Kl. 2 (Sedang) \& Tinggi & Kota Tasikmalaya \\
\hline
Kl. 2 (Sedang) \& Sedang & - \\
\hline
Kl. 2 (Sedang) \& Rendah & - \\
\hline
Kl. 3 (Tinggi) \& Sangat Tinggi & - \\
\hline
Kl. 3 (Tinggi) \& Tinggi & - \\
\hline
Kl. 3 (Tinggi) \& Sedang & Bandung Barat \\
\hline
Kl. 3 (Tinggi) \& Rendah & - \\
\hline
\end{tabular}
\end{adjustbox}
\end{table}